\section{Extensions}
\label{sec:ext}

Three extensions help interpret those answers: why an RDC may
saturate, when a classical reliability-growth curve emerges, and how
many steps are needed before $R(d)$ approaches $R_\infty$.

\subsection{RDC Shape}

An RDC is more useful when its shape has an explanation: a concave
curve suggests early steps resolve the easier cases while later ones
add less, and a convex window suggests a barrier such as waiting on a
tool response. The following condition captures the common
early-gains/saturation pattern.

\begin{assumption}[Stochastic monotonicity of $Q$]\label{A:monotone}
Fix an order $\preceq$ on the transient states. The matrix $Q$ is
\emph{stochastically monotone} if $Qf$ is non-increasing whenever
$f:\ST\to\mathbb{R}_{\ge0}$ is; equivalently, later rows of $Q$ place
stochastically more mass on later states in the
order~\cite{daley1968stochastic,karlin1968tp,keilsonKester1977}.
\end{assumption}

\begin{proposition}[RDC shape; Keilson--Kester~\cite{keilsonKester1977},
Karlin--Rinott~\cite{karlin1980tp2}]
\label{thm:shape}
Let $v=e_{s_0}^{\top}$ and $w=R_\oplus$. For
$R(d)=v(I-Q^d)NR_\oplus$,
\begin{equation}
\Delta R(d)=R(d+1)-R(d)=vQ^dw\ge0,
\label{eq:t4-1stdiff}
\end{equation}
and, for $d\ge1$,
\begin{equation}
\Delta^2R(d)=R(d+1)-2R(d)+R(d-1)
           =-vQ^{d-1}(I-Q)w.
\label{eq:t4-2nddiff}
\end{equation}
If $Qw\le w$ componentwise, then $R(d)$ is globally concave.
Assumption~\ref{A:monotone} together with $w$ non-increasing in the
state order is the structural regime in which that dominance is
expected; the concavity conclusion itself needs only $Qw\le w$.
\end{proposition}

\noindent\emph{Proof idea.}
The first difference is the probability of exiting to success exactly
after $d$ transient steps, so reliability never decreases with more
steps; since $vQ^{d-1}\ge0$, \eqref{eq:t4-2nddiff} makes $Qw\le w$
sufficient for a nonpositive second difference. A convex empirical
window instead signals an early barrier.

\subsection{NHPP Rare-Failure Limit}

The absorbing-chain view also explains when aggregate
reliability-growth curves emerge from step-level behavior---not as a
separate modeling story but as the rare-failure limiting case in which
a one-state chain reduces to Goel--Okumoto.

\begin{proposition}[NHPP rare-failure limit; Goel--Okumoto~\cite{goelokt}]
\label{thm:nhpp}
Consider a sequence of one-transient-state agent chains with per-step
success probability $\mu_n$, per-step failure probability
$\varepsilon_n$, and $Q_n=1-\mu_n-\varepsilon_n$. Suppose
$\mu_n\to0$, $\varepsilon_n\to0$, $d_n\mu_n\to\Lambda\in(0,\infty)$,
and $\varepsilon_n/\mu_n\to0$. For any fixed $c>0$ and $d=cd_n$, the
cumulative first-passage CDF converges to
\begin{equation}
R_n(d)\to 1-e^{-c\Lambda},
\label{eq:t5-GO}
\end{equation}
the Goel--Okumoto NHPP mean-value form $m(t)=a(1-e^{-bt})$ with
normalized $a=1$.
\end{proposition}

\noindent\emph{Proof idea.}
For the one-state chain,
$R_n(d)=\frac{\mu_n}{\mu_n+\varepsilon_n}
\left(1-(1-\mu_n-\varepsilon_n)^d\right)$. Under rare-failure
scaling the prefactor tends to one and the geometric term tends to
$e^{-c\Lambda}$. Outside this low-error regime, the exact closed form
of Proposition~\ref{thm:closed} should be used.

\subsection{Approach-Rate Bound}

Even when $R_\infty$ is known, a practitioner still needs a step
budget. The next result converts spectral decay of $Q$ into a
conservative horizon rule: how many more agent steps are needed before
the remaining reliability gap is at most $\delta$?

\begin{proposition}[Approach-rate bound; Horn--Johnson~\cite{hornjohnson2013},
Stewart~\cite{stewart1998}, Levin--Peres~\cite{levin2017markov}]
\label{thm:mixing}
Let $\rho=\rho(Q)<1$ and $N=(I-Q)^{-1}$. If
$Q=V\Lambda V^{-1}$ is diagonalizable, then
$R_\infty-R(d)\le\kappa_\infty(V)\rho^d\|NR_\oplus\|_\infty$
with $\kappa_\infty(V)=\|V\|_\infty\|V^{-1}\|_\infty$, so a
sufficient horizon for error at most $\delta$ is
\begin{equation}
d^\star_{\mathrm{diag}}(\delta)=
\max\!\left\{0,\left\lceil
\frac{\log(\kappa_\infty(V)\|NR_\oplus\|_\infty/\delta)}
     {-\log\rho}
\right\rceil\right\}.
\label{eq:t6-steps-diag}
\end{equation}
For non-diagonalizable $Q$, the same exponential term is multiplied
by a polynomial factor determined by the largest Jordan block.
\end{proposition}

\noindent\emph{Proof idea.}
The gap is $R_\infty-R(d)=e_{s_0}^{\top}Q^dNR_\oplus$, so bounding it
reduces to bounding $\|Q^d\|$; diagonalization gives
$\|Q^d\|_\infty\le\kappa_\infty(V)\rho^d$ and Jordan blocks add a
polynomial multiplier, so large $\rho(Q)$ or ill-conditioned
eigenvectors imply longer step budgets. Together these results connect
finite-horizon reliability, classical reliability growth, and
conservative step budgets within one fitted-chain view.

%% ================================================================
